\scanpage{99}
% Typescript inverted in the scan.
\mspage{-4-}
\uline{of regular classes of} $\mathrm{G}$.

In 7.4 an example of $\mathrm{N}$ is given: in case $\mathrm{G}$ is of type $\mathrm{\operatorname{SL}(r+1)}$
we obtain one of the classical normal forms under conjugacy. This special
case suggests the problem of extending the normal form $\mathrm{N}$ from regular
elements to arbitrary elements. In 7.1 it is shown that $\mathrm{N}$ is a closed
irreducible subset of $\mathrm{G}$, isomorphic as a variety to affine $\mathrm{r}$-space $\mathrm{V}$,
and in 7.9 (this is the main lemma concerning $\mathrm{N}$) that if $\mathrm{G}$ is simply
connected and $\mathrm{\chi_i}$ $\mathrm{(1\leq i\leq r)}$ denote the fundamental characters of $\mathrm{G}$
then the map $\mathrm{x\to(\chi_1(x),\chi_2(x),\ldots,\chi_r(x))}$ induces an isomorphism of $\mathrm{N}$
on $\mathrm{V}$. Then in \S8 the proof of 1.4 is given and simultaneously the following
important criterion for regularity is obtained.

1.5 Theorem. \uline{If} $\mathrm{G}$ \uline{is simply connected, the element} $\mathrm{x}$ \uline{is regular
if and only if the differentials} $\mathrm{d\chi_i}$ \uline{are independent at} $\mathrm{x}$.

At this point some words about recent work of B. Kostant are in order.
In [3] and [4] he has proved, among other thing, the analogues of our above
discussed results that are obtained by replacing the semisimple group $\mathrm{G}$ by
a semisimple Lie algebra $\mathrm{L}$ over the complex field (any algebraically
closed field of characteristic $\mathrm{0}$ will serve as well) and the characters
$\mathrm{\chi_i}$ of $\mathrm{G}$ by the basic polynomial invariants $\mathrm{u_i}$ of $\mathrm{L}$. The $\mathrm{\chi_i}$ turn
out to be considerably more tractable than the $\mathrm{u_i}$. Thus the proofs for $\mathrm{G}$
with no restriction on the characteristic are simpler than those for $\mathrm{L}$ in
characteristic $\mathrm{0}$. Assuming both $\mathrm{G}$ and $\mathrm{L}$ are in characteristic $\mathrm{0}$,
substantial parts of 1.1, 1.2, and 1.3 can be derived from their analogues
% archive sheet number: 98

